\chapter{电力电子器件}

就像晶体管和集成电路是模拟、数字电子电路的基础一样，\textbf{电力电子器件是电力电子电路的基础}。本章先概述器件的概念、特点与分类，再按“不可控器件$\rightarrow$半控型器件$\rightarrow$典型全控型器件$\rightarrow$其他新型器件”的顺序，逐个介绍其工作原理、基本特性、主要参数以及选择和使用中应注意的问题。器件应用中共性的驱动、保护、串并联问题留到第 9 章。

% ============================================================
\section{电力电子器件概述}

\subsection{电力电子器件的概念和特征}

在电气设备或电力系统中，\textbf{直接承担电能的变换或控制任务的电路称为主电路}（Main Power Circuit）。\textbf{电力电子器件是指可直接用于处理电能的主电路中，实现电能的变换或控制的电子器件}。

广义上电力电子器件也可分为电真空器件和半导体器件两类。但自 20 世纪 50 年代以来，除了在频率很高（如微波）的大功率高频电源中还在使用真空管外，基于半导体材料的电力电子器件已逐步取代汞弧整流器、闸流管等电真空器件，成为电能变换和控制领域的绝对主力，因此电力电子器件也往往\textbf{专指电力半导体器件}，主要材料仍是硅。

同处理信息的电子器件相比，电力电子器件具有如下特征：

\begin{enumerate}
  \item \textbf{所能处理电功率的大小（即承受电压和电流的能力）是其最重要的参数}。其处理电功率的能力小至毫瓦级、大至兆瓦级，一般都远大于处理信息的电子器件。
  \item 因为处理的电功率较大，为了减小本身的损耗、提高效率，电力电子器件一般都\textbf{工作在开关状态}。导通时（通态）阻抗很小、接近于短路，管压降接近于零，而电流由外电路决定；阻断时（断态）阻抗很大、接近于断路，电流几乎为零，而管子两端电压由外电路决定——就像普通晶体管的饱和与截止状态一样。因而器件的\textbf{动态特性（开关特性）和参数}也是其特性很重要的方面，有时甚至上升为第一位。
  \item 实际应用中电力电子器件往往需要由信息电子电路来控制。由于器件处理的电功率较大，普通信息电子电路的信号一般不能直接控制其导通或关断，需要中间电路对信号加以放大，这就是\textbf{驱动电路}。
  \item 尽管工作在开关状态，器件自身的功率损耗通常仍远大于信息电子器件，因而封装上讲究散热设计，工作时一般还需安装散热器。损耗包括\textbf{通态损耗}（通态压降 $\times$ 通态电流）、\textbf{断态损耗}（断态漏电流 $\times$ 断态电压）、\textbf{开关损耗}（开通损耗与关断损耗）。断态漏电流极其微小，故\textbf{通态损耗是主要成因}；当开关频率较高时，开关损耗会增大而可能成为主要因素。某些器件驱动电路注入的功率也是发热原因之一。
\end{enumerate}

\corepoint{把一个电力电子器件、或者外特性像一个开关的几个电力电子器件的组合，称为\textbf{电力电子开关}（电力半导体开关）；做电路分析时为简单起见往往用\textbf{理想开关}代替。}

\subsection{应用电力电子器件的系统组成}

电力电子器件在实际应用中，一般由\textbf{控制电路}、\textbf{驱动电路}和以电力电子器件为核心的\textbf{主电路}组成一个系统：由信息电子电路组成的控制电路按系统工作要求形成控制信号，通过驱动电路去控制主电路中器件的导通或关断，完成整个系统的功能。因此从宏观角度讲，电力电子电路也被称为\textbf{电力电子系统}。有的系统还需要\textbf{检测电路}采样主电路或现场信号；广义上把检测电路和驱动电路这些主电路以外的电路都归为控制电路，从而粗略地说，\textbf{电力电子系统由主电路和控制电路组成}。

主电路中的电压电流一般较大，而控制电路的元器件只能承受较小的电压电流，因此在主电路与控制电路连接的路径上（如驱动电路与主电路的连接处、驱动电路与控制信号的连接处、主电路与检测电路的连接处）一般需要\textbf{电气隔离}，通过光、磁等手段传递信号。此外，主电路中往往有电压电流过冲，而电力电子器件一般比主电路中普通元件昂贵、承受过电压过电流的能力却差一些，因此在主电路和控制电路中附加\textbf{保护电路}也往往非常必要。

电力电子器件一般都有三个端子：其中两个连接在主电路中、流通主电路电流，第三个称为\textbf{控制端（控制极）}。器件的导通或关断通过在控制端和一个主电路端子之间施加一定信号来控制，这个主电路端子是驱动电路和主电路的公共端，一般是主电路电流\textbf{流出}器件的那个端子。

\subsection{电力电子器件的分类}

\begin{table}[htbp]
\centering
\caption{电力电子器件的分类（按控制程度、驱动信号性质、载流子）}
\label{tab:2-1}
\begin{tabular}{p{2.7cm}p{4.8cm}p{6.5cm}}
\toprule
\textbf{分类依据} & \textbf{类别} & \textbf{代表器件与说明} \\
\midrule
\multirow{3}{2.4cm}{控制信号对通断的控制程度} & 不可控器件 & 电力二极管：不能用控制信号控制通断，因而也不需要驱动电路，只有两个端子；导通和关断完全由其在主电路中承受的电压电流决定 \\
 & 半控型器件 & 晶闸管及其大部分派生器件：控制信号可以控制其导通而\textbf{不能}控制其关断，关断完全由主电路承受的电压电流决定 \\
 & 全控型器件（自关断器件） & 通过控制信号\textbf{既可控制导通又可控制关断}；最常用的是 IGBT 和电力 MOSFET，此外还有 GTO、GTR 等 \\
\midrule
\multirow{2}{2.4cm}{驱动信号的性质} & 电流驱动型 & 通过从控制端注入或抽出电流来实现通断控制（晶闸管、GTO、GTR） \\
 & 电压驱动型（场控器件） & 仅通过在控制端和公共端之间施加一定电压信号即可控制通断；实质是用控制端电压在两个主电路端子间产生可控电场来改变电流，故又称场效应器件（电力 MOSFET、IGBT） \\
\midrule
\multirow{2}{2.4cm}{有效信号的波形} & 脉冲触发型 & 施加一个电压或电流\textbf{脉冲}即可开通或关断，此后只要主电路条件不变就能维持状态，不必继续施加控制信号（晶闸管、GTO） \\
 & 电平控制型 & 必须\textbf{持续}施加一定电平的信号才能开通并维持导通（或关断并维持阻断）（电力 MOSFET、IGBT） \\
\midrule
\multirow{3}{2.4cm}{导电载流子} & 单极型（多子器件） & 只有一种载流子参与导电：电力 MOSFET、SIT \\
 & 双极型（少子器件） & 电子和空穴两种载流子都参与导电：电力二极管、晶闸管、GTO、GTR \\
 & 复合型（混合型） & 由单极型与双极型器件集成混合而成：IGBT、MCT \\
\bottomrule
\end{tabular}
\end{table}

\subsection{本章内容和学习要点}

学习和选用电力电子器件时，最重要的是\textbf{掌握其基本特性}；此外应了解各国、各厂家对各种器件具体型号的命名方法，特别是要了解每种器件各个主要参数和特性曲线的意义。\textbf{掌握器件的型号命名法，以及其参数和特性曲线的使用方法，是在实际中正确应用电力电子器件的两个基本要求}。

许多电力电子器件都有其相对应的用于处理信息的电子器件（电力二极管、GTR、电力 MOSFET 分别与信息电子中的二极管、双极型晶体管、场效应晶体管对应），从半导体物理结构和工作原理上讲基本相同；但为了承受高电压和大电流，电力电子器件又具有与其信息电子对应者不同之处——这些不同之处，正是本章要抓的重点。

% ============================================================
\section{不可控器件——电力二极管}

电力二极管（Power Diode）自 20 世纪 50 年代初期就获得应用，当时也称为半导体整流器（SR），并已开始逐步取代汞弧整流器。虽然是不可控器件，但其结构和原理简单、工作可靠，所以直到现在仍大量应用于许多电气设备中；在采用全控型器件的电路中电力二极管往往不可缺少，特别是开通和关断速度很快的\textbf{快恢复二极管}和\textbf{肖特基二极管}具有不可替代的地位。

\subsection{PN 结与电力二极管的工作原理}

电力二极管实际上是由一个面积较大的 PN 结、两端引线及封装组成的。N 型半导体和 P 型半导体结合后构成 PN 结：交界处电子和空穴的浓度差别造成多子向另一区移动的\textbf{扩散运动}，在界面两侧留下带正、负电荷但不能任意移动的杂质离子，称为\textbf{空间电荷}；空间电荷建立的电场称为\textbf{内电场（自建电场）}，其方向阻止扩散运动，另一方面又吸引对方区内的少子向本区运动，这就是\textbf{漂移运动}。扩散与漂移相互联系又相互矛盾，最终达到动态平衡，形成稳定的\textbf{空间电荷区}（也称耗尽层、阻挡层或势垒区）。

\begin{itemize}
  \item \textbf{正向偏置}（外加电压正端接 P 区、负端接 N 区）：外加电场与自建电场方向\textbf{相反}，多子的扩散运动大于少子的漂移运动，空间电荷区变窄，外电路形成自 P 区流入、从 N 区流出的\textbf{正向电流} $I_F$——这是 PN 结的正向导通状态。
  \item \textbf{反向偏置}：外加电场与自建电场方向\textbf{相同}，少子的漂移运动大于多子的扩散运动，空间电荷区变宽，外电路形成自 N 区流入、从 P 区流出的\textbf{反向电流} $I_R$。少子浓度很小，温度一定时其数值趋于恒定，称为\textbf{反向饱和电流}，一般仅微安数量级，故反向偏置的 PN 结表现为高阻态——反向截止状态。
\end{itemize}

\textbf{PN 结的单向导电性}就是二极管的基本原理。为了建立承受高电压和大电流的能力，电力二极管的物理结构与信息电子二极管有如下不同：

\begin{enumerate}
  \item \textbf{垂直导电结构}：电流在硅片内流动的总体方向与硅片表面\textbf{垂直}（信息电子二极管一般是横向导电结构），使硅片中通过电流的有效面积增大，可显著提高通流能力。
  \item \textbf{P-i-N 结构}：在 P 区和 N 区之间多了一层\textbf{低掺杂 N 区}（记为 N$^-$），也称为\textbf{漂移区}（i 区）。低掺杂 N 区掺杂浓度低而接近本征半导体，可以承受很高的电压而不致被击穿——\textbf{低掺杂 N 区越厚，能承受的反向电压越高}。
  \item \textbf{电导调制效应}：低掺杂 N 区高电阻率对正向导通不利，这个矛盾通过电导调制效应解决。正向电流较小时，二极管电阻主要是作为基片的低掺杂 N 区的欧姆电阻，阻值较高且为常量，管压降随正向电流上升而增加；正向电流较大时，由 P 区注入并积累在低掺杂 N 区的少子空穴浓度将很大，为维持电中性，其多子浓度也相应大幅增加，使其电阻率明显下降（电导率大大增加）——这就是\textbf{电导调制效应}。它使电力二极管在正向电流较大时压降仍然很低，所以正向偏置的电力二极管表现为\textbf{低阻态}。
\end{enumerate}

PN 结具有一定的反向耐压能力，但当反向电压过大时反向电流将急剧增大，破坏反向偏置为截止的工作状态，这称为\textbf{反向击穿}，按机理分为\textbf{雪崩击穿}和\textbf{齐纳击穿}两种形式。反向击穿发生时，只要外电路采取措施将反向电流限制在一定范围内，则反向电压降低后 PN 结仍可恢复原来的状态；若反向电流未被限制住，使反向电流与反向电压的乘积超过 PN 结允许的耗散功率，就会因热量散发不出去导致结温上升直至过热烧毁，这就是\textbf{热击穿}。

PN 结中的电荷量随外加电压而变化，呈现电容效应，称为\textbf{结电容} $C_J$（微分电容），分为\textbf{势垒电容} $C_B$ 与\textbf{扩散电容} $C_D$：势垒电容只在外加电压变化时才起作用，其大小与 PN 结截面积成正比、与阻挡层厚度成反比；扩散电容仅在正向偏置时起作用。正向电压较低时以势垒电容为主，较高时扩散电容成为主要成分。结电容影响 PN 结的工作频率，高速开关状态下可能使其单向导电性变差甚至不能工作。

\subsection{电力二极管的基本特性}

\subsubsection*{1. 静态特性}

静态特性主要指\textbf{伏安特性}（图~\ref{fig:2-1}）。当承受的正向电压大到一定值（\textbf{门槛电压} $U_{TO}$）时，正向电流才开始明显增加，处于稳定导通状态；与正向电流对应的器件两端电压 $U_F$ 即其\textbf{正向电压降}。承受反向电压时，只有少子引起的微小而数值恒定的反向漏电流。

\begin{figure}[htbp]
\centering
\begin{tikzpicture}[font=\footnotesize,scale=0.95]
  \draw[->] (-4.2,0) -- (4.4,0) node[below] {$u$};
  \draw[->] (0,-3.1) -- (0,3.4) node[left] {$i$};
  \draw[yhbAccent,thick] (-4.0,-0.12) -- (-2.6,-0.12);
  \draw[yhbAccent,thick] (-2.6,-0.12) .. controls (-2.3,-0.9) and (-2.1,-2.6) .. (-1.9,-3.0);
  \draw[yhbAccent,thick] (0,0) -- (1.9,0.02);
  \draw[yhbAccent,thick] (1.9,0.02) .. controls (2.05,0.5) and (2.3,1.6) .. (2.6,3.0);
  \node[anchor=west] at (2.65,3.0) {$U_{TO}$};
  \node[anchor=west,text=textgray] at (-4.05,-0.45) {反向漏电流};
  \node[anchor=south,text=textgray] at (2.0,0.15) {门槛电压};
  \node[anchor=north] at (-2.4,0.1) {反向击穿};
  \draw[->,thin,textgray] (-1.2,-0.6) -- (-1.85,-0.9);
\end{tikzpicture}
\caption{电力二极管的伏安特性}
\label{fig:2-1}
\end{figure}

\subsubsection*{2. 动态特性}

因为结电容的存在，电力二极管在零偏置、正向偏置和反向偏置三种状态之间转换时必然经历一个过渡过程；过渡过程中 PN 结的一些区域需要一定时间调整带电状态，其电压—电流特性不能用伏安特性描述而是随时间变化的——这就是\textbf{动态特性}，并且往往专指反映通态与断态之间转换过程的\textbf{开关特性}。

\corepoint{{\bfseries 反向恢复过程}：原处于正向导通状态的电力二极管，外加电压突然由正向变为反向时，并\textbf{不能立即关断}，而需经过一段短暂时间才能重新获得反向阻断能力。关断之前有较大的反向电流出现，并伴随明显的反向电压过冲。原因是正向导通时在 PN 结两侧储存的大量少子需要被清除掉才能达到反向偏置稳态。}

设 $t=0$ 时刻外加电压突然由正向变为反向，正向电流在此反向电压作用下开始下降，下降的速率由主电路电感决定；到 $t=t_d$（延迟时间）时电流过零并开始反向增大，到 $t=t_f$ 时刻反向电流达到最大值 $I_{RP}$ 后开始减小，$t_f$ 称为电流下降时间。$t_d$ 与 $t_f$ 之和即\textbf{反向恢复时间} $t_{rr}$。其中 $t_d$ 内电荷存储效应起主要作用，$t_f$ 内则主要取决于结电容的充放电。

\keyformula{t_{rr}=t_d+t_f,\qquad Q_{rr}\approx \tfrac{1}{2}I_{RP}\,t_{rr}}

按反向恢复特性，二极管分为\textbf{软恢复特性}（反向电流下降较缓慢）与\textbf{硬恢复特性}（下降很陡、易引起较大电压过冲）两类。由零偏置转为正向偏置时也有开通过程（正向恢复过程），伴随电压过冲，但一般比反向恢复过程短得多、影响小。

\subsection{电力二极管的主要参数}

\begin{table}[htbp]
\centering
\caption{电力二极管的主要参数}
\label{tab:2-2}
\begin{tabular}{p{3.6cm}p{11.4cm}}
\toprule
\textbf{参数} & \textbf{含义与用法} \\
\midrule
正向平均电流 $I_{F(AV)}$ & 在规定的管壳温度和散热条件下，允许流过的\textbf{最大工频正弦半波电流的平均值}。定额按发热效应（即有效值相等）等效，正弦半波电流的\textbf{平均值:有效值 $=1:1.57$}，故实际波形电流有效值 $I$ 对应 $I_{F(AV)}=I/1.57$，选用时再留 1.5$\sim$2 倍裕量 \\
正向压降 $U_F$ & 在规定温度下流过规定电流时器件两端的正向电压降 \\
反向重复峰值电压 $U_{RRM}$ & 允许重复施加的反向最高峰值电压，通常为反向不重复峰值电压 $U_{RSM}$ 的 $90\%$；使用时取正常工作峰值电压的 2$\sim$3 倍 \\
反向漏电流 $I_R$ & 在额定结温、施加反向重复峰值电压时的反向电流 \\
最高工作结温 $T_{jm}$ & 通常为 $125\sim175\,^\circ$C，取决于器件封装与散热条件 \\
浪涌电流 $I_{FSM}$ & 所能承受的最大的连续一个或几个工频周期的过电流 \\
反向恢复时间 $t_{rr}$ & 关断过程中重新获得反向阻断能力所需时间 \\
\bottomrule
\end{tabular}
\end{table}

\subsection{电力二极管的主要类型}

\begin{itemize}
  \item \textbf{普通二极管}（Standard Recovery Diode）：又称整流二极管，多用于开关频率不高（1 kHz 以下）的整流电路，其反向恢复时间较长、在 $5\,\mu$s 以上，正向电流定额和反向电压定额却可以达到很高。
  \item \textbf{快恢复二极管}（Fast Recovery Diode，FRD）：反向恢复时间很短（$5\,\mu$s 以下），正向压降较低（约 0.9 V），多用于开关频率较高的电路，如与全控型器件配合的续流二极管。从性能上分为\textbf{快速恢复}和\textbf{超快恢复}两种。
  \item \textbf{肖特基二极管}（Schottky Barrier Diode，SBD）：以金属—半导体结代替 PN 结，反向恢复时间极短（几十纳秒），正向压降很低（0.4$\sim$0.6 V）；但反向耐压较低（一般不超 100 V 左右）、反向漏电流较大，多用于低电压、高频场合。
\end{itemize}

% ============================================================
\section{半控型器件——晶闸管}

晶闸管（Thyristor）是晶体闸流管的简称，又称可控硅整流器（SCR），是目前\textbf{容量最大、应用最广泛}的电力电子器件之一。

\subsection{晶闸管的结构与工作原理}

晶闸管是大功率四层（P1-N1-P2-N2）三端（阳极 A、阴极 K、门极 G）器件，内部有三个 PN 结 J$_1$、J$_2$、J$_3$。从外形看主要有螺栓形、平板形等封装。

\begin{center}
\begin{circuitikz}[font=\footnotesize,scale=0.9]
  \draw (0,0) node[left] {A} to[short,o-] (1,0) to[thyristor] (3.4,0) to[short,-o] (4.4,0) node[right] {K};
  \draw (2.55,-0.95) -- (1.95,-1.55) node[below] {G};
\end{circuitikz}
\end{center}

\textbf{工作原理（双晶体管等效）}：把晶闸管看成 P1-N1-P2 构成的 PNP 晶体管 V$_1$ 与 N1-P2-N2 构成的 NPN 晶体管 V$_2$ 的互联——V$_1$ 的集电极即 V$_2$ 的基极，V$_2$ 的集电极即 V$_1$ 的基极。设两个晶体管的共基极电流放大系数分别为 $\alpha_1$、$\alpha_2$，则阳极电流

\keyformula{I_A=\frac{\alpha_2 I_G+\left(\alpha_1+\alpha_2\right)I_{CBO}}{1-\left(\alpha_1+\alpha_2\right)}}

\corepoint{晶体管的特性是：\textbf{低发射极电流下 $\alpha$ 很小，发射极电流增大时 $\alpha$ 迅速增大}。因此当门极注入电流 $I_G$ 后，两个晶体管相互正反馈，$\alpha_1+\alpha_2$ 迅速趋近于 1，$I_A$ 急剧增大，晶闸管\textbf{饱和导通}。导通后 $\alpha_1+\alpha_2\approx1$，即使门极电流消失，晶闸管仍维持导通——\textbf{门极失去控制作用}。要使晶闸管关断，只有设法使阳极电流减小到维持电流以下，使 $\alpha_1+\alpha_2<1$。}

由此得到晶闸管导通与关断的条件：

\begin{table}[htbp]
\centering
\caption{晶闸管的导通与关断条件}
\label{tab:2-3}
\begin{tabular}{p{2.2cm}p{12.8cm}}
\toprule
\textbf{状态} & \textbf{条件} \\
\midrule
由断态转通态 & （1） 阳极电位高于阴极（\textbf{承受正向阳极电压}）；（2） 门极承受\textbf{正向门极电压并有足够的门极电流} $I_G$。两者\textbf{缺一不可} \\
维持通态 & 阳极电流大于\textbf{维持电流} $I_H$；此时门极信号已不起作用 \\
由通态转断态 & 阳极电流减小到维持电流以下（或阳极电压反向）。方法是：主电路电流过零、外加反向电压、或电路阻抗增大使电流降到 $I_H$ 以下 \\
\bottomrule
\end{tabular}
\end{table}

\subsection{晶闸管的基本特性}

\subsubsection*{1. 静态特性}

\begin{itemize}
  \item \textbf{正向特性}：门极断路（$I_G=0$）时，逐渐增大正向阳极电压，只有很小的正向漏电流，晶闸管处于\textbf{正向阻断状态}；当正向电压增大到\textbf{正向转折电压} $U_{BO}$ 时，漏电流急剧增大，由断态转入通态（这种靠过电压导通的方式叫“硬开通”，会损坏器件，正常不用）。门极施加电流 $I_G$ 后，正向转折电压随 $I_G$ 增大而降低；晶闸管导通后，通态压降很小（约 1 V 左右），流过很大的阳极电流。
  \item \textbf{反向特性}：晶闸管承受反向电压时，只有极小的反向漏电流，处于\textbf{反向阻断状态}；反向电压增大到一定程度时反向电流急剧增大，即\textbf{反向击穿}，器件损坏。
\end{itemize}

\begin{center}
\begin{tikzpicture}[font=\footnotesize,scale=0.95]
  \draw[->] (-4.0,0) -- (4.6,0) node[below] {$u$};
  \draw[->] (0,-2.9) -- (0,3.5) node[left] {$i$};
  % 正向阻断段
  \draw[yhbAccent,thick] (0,0) -- (2.75,0.05);
  \draw[yhbAccent,thick] (2.75,0.05) .. controls (2.95,0.7) and (3.0,1.6) .. (3.05,3.1);
  % 通态段
  \draw[yhbAccent,thick] (2.75,0.05) .. controls (2.2,1.1) and (1.6,2.2) .. (1.45,3.1);
  % 反向
  \draw[yhbAccent,thick] (-3.6,-0.08) -- (-2.9,-0.1);
  \draw[yhbAccent,thick] (-2.9,-0.1) .. controls (-2.75,-0.7) and (-2.65,-1.9) .. (-2.6,-2.7);
  \node[anchor=west] at (3.1,3.05) {$U_{BO}$};
  \node[anchor=east] at (1.4,3.05) {通态};
  \node[anchor=west] at (2.78,1.9) {\small $I_G$ 增大};
  \node[anchor=north,textgray] at (-2.55,-0.15) {\small 反向击穿};
  \node[anchor=south,textgray] at (2.2,0.2) {\small 正向阻断};
\end{tikzpicture}
\end{center}

\subsubsection*{2. 动态特性}

\begin{itemize}
  \item \textbf{开通过程}：门极加触发脉冲后，阳极电流滞后一段\textbf{延迟时间} $t_d$（约零点几到 1.5 $\mu$s）才明显上升，随后是\textbf{上升时间} $t_r$（约零点几到 3 $\mu$s）；\textbf{开通时间} $t_{gt}=t_d+t_r$。开通时须注意\textbf{通态电流临界上升率 di/dt}：上升太快则刚导通时电流集中在门极附近的小区域，局部过热而损坏器件（须串电感或用强触发）。
  \item \textbf{关断过程}：先把阳极电流降到零并施加反压，器件内部存储电荷需要清除，存在\textbf{反向阻断恢复时间} $t_{rr}$；之后还要经过\textbf{正向阻断恢复时间} $t_{gr}$ 才恢复正向阻断能力。两者之和即\textbf{关断时间}

\keyformula{t_q=t_{rr}+t_{gr}}

  在正向阻断恢复时间内若重新施加正向电压，晶闸管会重新正向导通而不是由门极控制导通，故实际应用中须施加足够长时间的反压使其充分恢复阻断能力。
\end{itemize}

\corepoint{{\bfseries 断态电压临界上升率 du/dt}：在额定结温和门极开路情况下，不导致晶闸管从断态到通态转换的外加电压最大上升率。正向电压上升过快时，处于阻断状态的 J$_2$ 结（相当于一个电容）会有充电\textbf{位移电流}流过，流经 J$_3$ 时起类似门极触发电流的作用，\textbf{电压上升率过大会使晶闸管误导通}。使用中实际电压上升率必须低于此临界值（可加 RC 吸收电路）。}

\subsection{晶闸管的主要参数}

普通晶闸管在反向稳态下一定是阻断状态；而正向工作时既可能导通也可能阻断，因此提到参数时，\textbf{断态和通态都是为了区分正向的不同状态，“正向”二字可省去}。各项主要参数往往与结温相联系。

\subsubsection*{1. 电压定额}

\begin{itemize}
  \item \textbf{断态重复峰值电压} $U_{DRM}$：门极断路、结温为额定值时，允许重复加在器件上的正向峰值电压（国标规定重复频率 50 Hz、每次持续时间不超过 10 ms）。规定 $U_{DRM}$ 为\textbf{断态不重复峰值电压}（断态最大瞬时电压）的 $90\%$，后者应低于正向转折电压。
  \item \textbf{反向重复峰值电压} $U_{RRM}$：同样条件下允许重复加在器件上的反向峰值电压，为反向不重复峰值电压的 $90\%$，后者应低于反向击穿电压。
  \item 通常取 $U_{DRM}$ 与 $U_{RRM}$ 中\textbf{较小}的标称值作为该器件的\textbf{额定电压}；选用时额定电压要留裕量，一般取为正常工作时器件承受峰值电压的\textbf{2$\sim$3 倍}。
  \item \textbf{通态（峰值）电压} $U_{TM}$：通以某一规定倍数额定通态平均电流时的通态峰值电压。
\end{itemize}

\subsubsection*{2. 电流定额}

\begin{itemize}
  \item \textbf{通态平均电流} $I_{T(AV)}$：在环境温度 40 $^\circ$C 和规定冷却状态下，稳定结温不超过额定结温时所允许流过的\textbf{最大工频正弦半波电流的平均值}，也是标称其额定电流的参数。与电力二极管一样，该参数按发热效应（有效值相等）等效，正弦半波\textbf{平均值:有效值 $=1:1.57$}，故先由实际波形有效值算出 $I_{T(AV)}=I/1.57$，再取\textbf{1.5$\sim$2 倍}裕量。
  \item \textbf{维持电流} $I_H$：使晶闸管维持导通所必需的最小电流，一般为几十到几百毫安；结温越高 $I_H$ 越小。
  \item \textbf{擎住电流} $I_L$：晶闸管刚从断态转入通态并\textbf{移除触发信号后}能维持导通所需的最小电流。对同一器件通常 $I_L\approx(2\sim4)I_H$。
  \item \textbf{浪涌电流} $I_{FSM}$：电路异常情况引起的、使结温超过额定结温的不重复性最大正向过载电流，可作为设计保护电路的依据。
\end{itemize}

\corepoint{选型步骤记牢三步：（1） 由实际波形求电流\textbf{有效值} $I$；（2） 折算到正弦半波平均值 $I_{T(AV)}=I/1.57$；（3） 乘 1.5$\sim$2 倍裕量后取标准档。电压则按承受的\textbf{峰值}乘 2$\sim$3 倍。}

\subsubsection*{3. 动态参数}

除开通时间 $t_{gt}$ 和关断时间 $t_q$ 外，还有\textbf{断态电压临界上升率 du/dt} 与\textbf{通态电流临界上升率 di/dt}。普通晶闸管 $t_q$ 约为数百微秒，快速晶闸管为数十微秒，高频晶闸管为 10 $\mu$s 左右。

\subsection{晶闸管的派生器件}

\begin{itemize}
  \item \textbf{快速晶闸管}（FST）：包括所有专为快速应用而设计的晶闸管，有常规快速晶闸管和高频晶闸管，可分别应用于 400 Hz 和 10 kHz 以上的斩波或逆变电路。其开关时间以及 du/dt、di/dt 耐量都明显改善；不足是电压和电流定额都不易做高，且因工作频率较高，选择通态平均电流时\textbf{不能忽略开关损耗的发热效应}。
  \item \textbf{双向晶闸管}（TRIAC）：可认为是一对反并联普通晶闸管的集成，有两个主电极 T$_1$、T$_2$ 和一个门极 G。门极使器件在主电极的正反两个方向均可触发导通，因此在第 I、第 III 象限有对称的伏安特性。与一对反并联晶闸管相比更经济、控制电路简单，广泛用于\textbf{交流调压电路}、固态继电器（SSR）和交流电动机调速。
  \item \textbf{逆导晶闸管}（RCT）：将晶闸管和二极管反并联集成在同一管芯上，其正向压降小、关断时间短、高温特性好、额定结温高；用于不需要阻断反向电压的电路中（如斩波、逆变中的续流）。
  \item \textbf{光控晶闸管}（LTT）：用光信号代替电信号触发，门极触发改为光触发，具有主电路与控制电路间天然电气隔离、抗干扰能力强的优点，多用于高压直流输电等超高压大容量场合。
\end{itemize}

% ============================================================
\section{典型全控型器件}

20 世纪 70 年代后期，以 GTO、GTR 和电力 MOSFET 为代表的全控型器件迅速发展，使电力电子技术的面貌焕然一新。本节按 GTO $\rightarrow$ GTR $\rightarrow$ 电力 MOSFET $\rightarrow$ IGBT 的顺序介绍。

\subsection{门极可关断晶闸管（GTO）}

GTO（Gate Turn-Off Thyristor）是晶闸管的一种派生器件，但可以通过在门极施加负的脉冲电流使其关断，因而属于\textbf{全控型器件}。其电压、电流容量较大，与普通晶闸管接近，在兆瓦级以上的大功率场合仍有较多应用。

\begin{itemize}
  \item \textbf{结构特点}：与普通晶闸管的 PNPN 四层结构相同，但\textbf{多元集成结构}——一个 GTO 由数百甚至上千个小 GTO 单元并联而成，阴极条做得细长，以便门极电流能迅速均匀地流过整个阴极面积。门极为横向结构、阴极条窄而长，是它能实现门极关断的结构基础。
  \item \textbf{工作原理}：GTO 导通时 $\alpha_1+\alpha_2$ 比 1 略大但很接近 1（普通晶闸管则远大于 1 进入饱和），处于\textbf{临界饱和}状态，这为门极关断创造了条件。关断时门极加负电流 $I_{GM}$ 抽出存储电荷，使 $\alpha_1+\alpha_2<1$，阳极电流下降并关断。
  \item \textbf{电流关断增益} $\beta_{off}$：最大可关断阳极电流 $I_{ATO}$ 与门极负电流最大值 $I_{GM}$ 之比（注：此处 $I_{ATO}$ 即手册中的 $I_{ATM}$，习惯写作 $I_{ATO}$）
\end{itemize}

\keyformula{\beta_{off}=\frac{I_{ATO}}{I_{GM}}}

$\beta_{off}$ 一般很小，只有 5 左右，这是 GTO 的一个主要缺点——需要很大的门极负脉冲电流才能关断，驱动功率大。

\subsection{电力晶体管（GTR）}

电力晶体管（Giant Transistor，GTR）是一种双极型大功率晶体管，通常采用\textbf{共发射极接法}，由基极电流控制集电极电流。它是双极型电流驱动器件，有电导调制效应，通态压降较低。

\begin{itemize}
  \item \textbf{输出特性}：明显分为截止区、放大区和饱和区。作为开关使用时，工作在\textbf{截止区（断态）}和\textbf{饱和区（通态）}，放大区只是过渡。饱和时管压降很低（$U_{CE(sat)}$），截止时只有极小的漏电流。
  \item \textbf{主要参数}：$U_{CEO}$（$I_B=0$ 时集电极—发射极击穿电压）等一组击穿电压、集电极最大电流 $I_{CM}$、集电极最大耗散功率 $P_{CM}$、共射极电流放大系数 $h_{FE}$（开通时是 $h_{FE}$，关断时须注意存储时间）。
  \item \textbf{二次击穿}：GTR 特有的失效机理。当集电极电压较高、电流较大时，局部热点使电流更集中于该点，形成正反馈，导致器件在远低于额定功耗的条件下瞬间损坏。因此 GTR 有\textbf{安全工作区（SOA）}，分正向偏置 SOA 与反向偏置 SOA，使用时必须保证工作点落在 SOA 之内，并加\textbf{缓冲电路}把开关轨迹限制在安全区内。
  \item GTR 的开关时间在几微秒量级，比晶闸管短得多，但仍远长于电力 MOSFET；驱动电流大、驱动功率大是其主要缺点，现已逐步被 IGBT 取代。
\end{itemize}

\subsection{电力场效应晶体管（Power MOSFET）}

电力 MOSFET 是小功率 MOSFET 的功率化，属于\textbf{单极型电压驱动器件}，只靠\textbf{多子}导电，因而开关速度极快。

\begin{itemize}
  \item \textbf{结构}：采用\textbf{垂直导电结构}（VVMOS、VDMOS），有 N 沟道和 P 沟道两种，以 N 沟道（N-MOSFET）为主；内部\textbf{源极与衬底短接}，并内含反并联的体二极管。由于没有少子存储效应，其开关过程只是对输入电容（$C_{iss}$）的充放电。
  \item \textbf{静态特性}：输出特性分为\textbf{截止区、饱和区（恒流区）、可变电阻区}。在开关应用中工作在截止区与可变电阻区之间切换；导通时表现为一个\textbf{通态电阻} $R_{on}$（$R_{DS(on)}$），因此通态压降随电流线性增大，这是它在高压大电流场合通态损耗偏大的原因。
  \item \textbf{转移特性}：漏极电流 $I_D$ 与栅源电压 $U_{GS}$ 的关系，$U_{GS}$ 超过\textbf{开启电压} $U_{GS(th)}$ 后才出现 $I_D$；跨导 $g_{fs}=\Delta I_D/\Delta U_{GS}$ 表征其放大能力（作开关用时关心的是 $R_{on}$ 与栅极电荷）。
  \item \textbf{动态特性}：开通、关断过程即栅极电容充放电过程，开关时间在几十到几百纳秒量级；开关损耗与驱动电阻、栅极电荷有关。
\end{itemize}

\corepoint{电力 MOSFET 的三大优点：（1） \textbf{驱动功率小}（电压驱动、输入阻抗高、静态几乎不耗驱动功率）；（2） \textbf{开关速度极快}（无少子存储效应）；（3） \textbf{无二次击穿}，且 $R_{on}$ 具有\textbf{正温度系数}，多管并联时能自动均流。主要缺点：\textbf{通态压降大}（耐压越高 $R_{on}$ 越大，故耐压不易做高），一般用于 1000 V 以下、高频场合。}

\subsection{绝缘栅双极晶体管（IGBT）}

GTR 和 GTO 是双极型电流驱动器件，电流驱动、驱动功率大、开关速度较低；电力 MOSFET 是单极型电压驱动器件，驱动功率小、速度快但通态压降大。IGBT 把二者的优点结合起来，成为\textbf{现代电力电子技术的主导器件}。

\begin{itemize}
  \item \textbf{结构}：由沟道型 VDMOSFET 与双极型晶体管组合而成的复合器件，相当于在 MOSFET 的漏极侧再加一层 P$^+$ 区，形成四层结构。对外仍是三个电极：\textbf{栅极 G、集电极 C、发射极 E}（对应 MOSFET 的 G、D、S）。
  \item \textbf{特性}：\textbf{电压驱动}（栅极输入阻抗高、驱动功率小），\textbf{通态压降小}（因为 P$^+$ 区向漂移区注入少子，产生电导调制效应）；转移特性与输出特性与 MOSFET 相似，输出特性同样有截止区、饱和区、可变电阻区，但饱和区（恒流区）较窄。
  \item \textbf{参数}：$U_{CES}$（集电极—发射极最高电压）、$I_C$（额定集电极电流）、$U_{CE(sat)}$（通态饱和压降）、$U_{GE(th)}$（开启电压，一般 3$\sim$6 V）、开关时间等。
  \item \textbf{掣住效应（latch-up）}：IGBT 内部天然存在一个寄生晶闸管（PNPN），当集电极电流过大或关断过快时寄生晶闸管可能被触发导通，栅极失去控制作用，导致器件损坏。这是 IGBT 使用中必须避免的失效模式，须限制最大集电极电流、必要时加负栅压关断。
\end{itemize}

\begin{table}[htbp]
\centering
\caption{几种典型电力电子器件的对比}
\label{tab:2-4}
\begin{tabular}{p{2.5cm}p{2.5cm}p{2.4cm}p{2.4cm}p{3.6cm}}
\toprule
\textbf{器件} & \textbf{驱动方式} & \textbf{控制程度} & \textbf{开关速度} & \textbf{特点} \\
\midrule
电力二极管 & 无需驱动 & 不可控 & 取决于类型 & 结构简单、可靠；快恢复/肖特基用于高频 \\
晶闸管 & 电流（脉冲） & 半控 & 慢（$t_q$ 数百 $\mu$s） & 容量最大、耐压最高；需强迫换流关断 \\
GTO & 电流（脉冲） & 全控 & 较慢 & 容量大、可关断；$\beta_{off}$ 小、驱动功率大 \\
GTR & 电流（持续） & 全控 & 中（几 $\mu$s） & 通态压降小；有二次击穿、驱动功率大 \\
电力 MOSFET & 电压 & 全控 & 快（ns 级） & 驱动功率小、$R_{on}$ 正温度系数；耐压不易高 \\
IGBT & 电压 & 全控 & 较快 & 兼具 MOSFET 驱动优点与 GTR 通态优点；须防掣住 \\
\bottomrule
\end{tabular}
\end{table}

% ============================================================
\section{其他新型电力电子器件}

\subsection{MOS 控制晶闸管（MCT）}

MCT（MOS Controlled Thyristor）是把 MOSFET 与晶闸管集成的复合器件，可在门极加信号控制其开通与关断，属于电压驱动型全控器件。其特点是：通态压降小、通流能力强、du/dt 与 di/dt 耐量高、开关速度快；缺点是制造工艺复杂、成本高。MCT 一度被看作很有前途的大功率器件，但实用化进展不如 IGBT。

\subsection{静电感应晶体管（SIT）}

SIT（Static Induction Transistor）是一种结型场效应器件，为\textbf{多子导电的单极型器件}，具有输出功率大、输入阻抗高、开关速度快、无二次击穿、抗辐射能力强等优点。缺点是通态电阻较大、驱动电路复杂，多用于高频大功率放大与感应加热电源。

\subsection{静电感应晶闸管（SITH）}

SITH（Static Induction Thyristor）又称场控晶闸管（FCTh），是在 SIT 基础上发展起来的，属\textbf{双极型}器件。其通态压降小、通流能力强、开关速度快、du/dt 与 di/dt 耐量高；缺点是需要较大的驱动功率、关断时存在尾部电流。

\subsection{集成门极换流晶闸管（IGCT）}

IGCT（Integrated Gate Commutated Thyristor）是 GTO 与 MOS 驱动电路集成而成的新型器件，实质是把门极驱动电路与 GTO 管芯集成在一起，实现“门极换流”。其特点是：容量大（可达数千安、数千伏）、通态压降小、开关频率比 GTO 高、无需缓冲电路、可靠性高，在大功率变频器、电力系统（如静止同步补偿器、直流输电）中已获得大量应用。

\subsection{基于宽禁带半导体材料的电力电子器件}

到目前为止，硅材料一直是电力电子器件的主要材料。硅的禁带宽度为 1.12 eV，而\textbf{宽禁带半导体材料}是指禁带宽度在 3.0 eV 及以上的半导体材料，典型的是\textbf{碳化硅（SiC）、氮化镓（GaN）}等。宽禁带材料具有临界击穿场强高、热导率高、饱和电子漂移速度高等优点，用其制成的器件\textbf{耐压更高、通态电阻更小、可在更高温度和更高频率下工作}，是电力电子器件的重要发展方向。目前 SiC 肖特基二极管、SiC MOSFET 已商品化，但成本、可靠性与驱动技术仍有很长的路要走。

% ============================================================
\section{功率集成电路与集成电力电子模块}

自 20 世纪 80 年代中后期开始，在电力电子器件研制和开发中的一个共同趋势是\textbf{集成化}，这就是电力电子集成电路（PIC）的概念。

\begin{itemize}
  \item \textbf{功率模块}（Power Module）：把若干个电力电子器件及必要的辅助元件封装成一个模块，给应用带来很大方便——结构紧凑、体积减小、可靠性提高。
  \item \textbf{集成电力电子模块}（IPEM）/ \textbf{智能功率模块}（IPM）：进一步把驱动、控制、保护（含过流、过热、欠压保护）电路与电力电子器件集成在一起，构成“功率 + 驱动 + 保护”的模块。用户只需给输入信号，模块自己完成保护动作，这在实际工程中极为常用。
  \item \textbf{单片集成}：把主电路器件与控制电路做在同一芯片上，难度最大，受制于\textbf{电压隔离}（主电路高压、控制电路低压）、\textbf{热隔离}（主电路发热严重）、\textbf{电磁干扰}（开关器件通断高压大电流且与控制电路同处一芯）等难题，目前功率都还较小。
  \item \textbf{混合集成}：把不同的单个芯片集成封装在一起，功率密度不如单片集成，但为解决上述难题提供了很大方便，\textbf{封装技术成为关键技术}，是当前电力电子集成技术发展的焦点。此外\textbf{系统集成}也是重要方面，对超大功率集成技术尤其如此。
\end{itemize}

\corepoint{本章复习主线：\textbf{“控制程度—驱动性质—载流子”三把尺子}给器件定位；\textbf{晶闸管}抓“门极只开不关 + 双晶体管正反馈 + 选型三步”；\textbf{全控型器件}抓“驱动功率与通态压降的折中”——GTR 通态好但驱动费劲、MOSFET 驱动省但通态差、IGBT 取二者之长，这就是第 5—8 章各种电路选用器件的依据。}
